\section{Solution method}
\label{sec:method}

\subsection{Formulation selection}

The direct implementation of Section~\ref{sec:model} is a mixed-integer program with one binary segment
variable for every demand, period, and reachable ordered pair of routers. This compact representation
is exact for the scalar MLU model and exposes flow conservation, segment limits, link loads, and, in
the two-period case, reconfiguration decisions directly to the solver. Its principal limitation is the
$O(|D||T||V|^2)$ number of routing variables. The implementation therefore estimates this number as
$|D||T||V|(|V|-1)$ and uses the compact formulation when the estimate is at most $100\,000$; larger
instances are solved by the path-based method described below.

\subsection{Single-period column generation}

For a period~$t$, let $\mathcal{P}_{dt}$ be the set of admissible segment paths for demand~$d$, and
let $R_{pat}$ denote the fraction of path~$p$ carried by physical link~$a$. A master variable
$z_{dtp}$ selects one complete path, so flow conservation and the segment limit are enforced by the
definition of a column. For a restricted subset $\widehat{\mathcal{P}}_{dt}\subseteq\mathcal{P}_{dt}$,
the LP master is
\begin{align}
\min\ & L, \nonumber\\
\text{s.t.}\ & \sum_{p\in\widehat{\mathcal{P}}_{dt}} z_{dtp}=1, && \forall d, \label{eq:convexity}\\
& \sum_{d\in D}\sum_{p\in\widehat{\mathcal{P}}_{dt}} v_{dt}\,R_{pat}\,z_{dtp}\le c_a L, && \forall a\in A_t, \label{eq:master-load}\\
& z_{dtp}\ge 0. \nonumber
\end{align}
The initial pool contains the direct source--destination segment for every demand. After solving the
restricted master, pricing uses the dual values of \eqref{eq:convexity}--\eqref{eq:master-load} to
find a minimum-reduced-cost path with at most $M$ segments. This bounded shortest-path problem is
solved exactly by dynamic programming on $M$ layers of the complete segment graph.

Let $\bar c^*_{dt}$ be the minimum reduced cost returned for demand~$d$. At every pricing round,
\[
\mathrm{LB}_t=z_{\mathrm{RMP}}+\sum_{d\in D}\min\{0,\bar c^*_{dt}\}
\]
is a valid Lagrangian lower bound for the complete path master. The dynamic program minimizes over
bounded walks, a superset of simple paths; therefore its value remains a valid lower estimate even if
a minimum walk repeats a router. Cycles are removed before a feasible column is inserted. If no
negative-reduced-cost path exists, the restricted LP value itself is the optimum of the complete LP
relaxation. The generated variables are then made binary and the restricted integer master is solved
to obtain a feasible routing and an upper bound.

\subsection{Joint two-period solution}

For the full experiment, column generation is first applied separately to periods~0 and~1. Each solve
receives 40\% of the total time limit and returns an initial path pool, an incumbent routing, and a
valid single-period lower bound. Because the budget constraint only restricts the joint feasible set,
$\max\{\mathrm{LB}_0,\mathrm{LB}_1\}$ is a valid initial lower bound for the full problem. The two
pools are then enriched with a small number of feasible waypoint paths and with routes that are valid
in both period graphs.

The pools are coupled in a joint master containing one path-selection constraint per demand and
period, the link-load constraints of both periods, and the absolute-change
linearization~\eqref{eq:change}--\eqref{eq:budget}. An initial binary joint master, warm-started from
a repaired budget-feasible pair of routings, provides a joint incumbent. The corresponding LP
relaxation is then used for joint pricing. The reduced cost of a path includes the dual contribution
of its period's link-load constraints, its demand-selection constraint, and the two absolute-change
rows associated with every segment already represented in the joint master. For a previously unseen
segment, assigning equal dual values to its two new change rows gives zero net reconfiguration
penalty; this is a dual-feasible completion and hence preserves the validity of the lower bound.

Pricing is performed for every demand in both periods by the same bounded dynamic program. If
$\bar c^*_{dt}$ is the resulting minimum reduced cost, the joint bound is
\[
\mathrm{LB}_{\mathrm{joint}}=z^{\mathrm{LP}}_{\mathrm{joint}}+\sum_{t\in\{0,1\}}\sum_{d\in D}\min\{0,\bar c^*_{dt}\}.
\]
Negative-reduced-cost paths are added and the joint LP is re-solved. The implementation certifies the
complete joint LP when the correction term is zero. Practical safeguards limit joint pricing to 40
rounds and $25\,000$ columns and stop it after six rounds without LP improvement; if a safeguard is
reached first, the corrected Lagrangian value remains a valid, possibly weaker, bound. The reported
global lower bound is
\[
\mathrm{LB}=\max\{\mathrm{LB}_0,\mathrm{LB}_1,\mathrm{LB}_{\mathrm{joint}}\}.
\]
Finally, the enlarged budget-coupled master is solved with binary path variables.

\subsection{Use of the remaining time}

Obtaining a valid relaxation bound does not imply that the best integer routing in the current column
pool attains that bound. When a positive gap and usable time remain, the implementation therefore
generates additional feasible columns using systematic one-waypoint paths, randomized two-waypoint
paths, and bounded paths obtained from randomized segment costs. It warm-starts the enlarged binary
master from the best validated routing and solves it again with a different random seed. For the
period-0 experiment, this diversification is applied to the single-period master. For the full
experiment, the independent period-level diversification is disabled after the initial pools have been
constructed; a candidate is instead added to every period pool in which it is feasible, and the
budget-coupled joint MIP is re-solved. The process stops when the gap closes, the time limit is
reached, or no new path is found and the integer solve terminates immediately. These heuristic rounds
can improve only the feasible upper bound; they do not modify the previously established relaxation
lower bound.

\subsection{Bounds and certification}

After optimization, every incumbent is checked independently for correct terminals, reachability,
segment count, repeated vertices, reconfiguration cost, and recomputed link loads. The checked maximum
load is the reported upper bound UB. The compact solver or the column-generation procedure supplies
the lower bound LB, and the reported gap is
\[
\max\left\{0,\,1-\frac{\mathrm{LB}}{\mathrm{UB}}\right\}.
\]
A run is labelled \emph{MLU-optimal} when this relative gap is at most $10^{-6}$, the numerical
tolerance adopted by the implementation. This certification concerns the scalar
objective~\eqref{eq:mlu}; it does not certify the remaining ranks of the challenge's lexicographic
objective.
